\section{Protocol and essential audit notes}
\label{app:protocol}

\paragraph{Data and scoring.}
The preparation pipeline samples 600,000 BlueGene/L log lines by reservoir
sampling (seed 42) and splits them temporally into 400,000 training and
200,000 test records. Evaluation uses subsets of that test split. References
come from the repository's parser, which supports Drain~\cite{drain} and
a regular-expression fallback; they are not independent semantic labels.
The anomaly rule marks a record positive if its native alert tag is
non-normal or its severity is ERROR, FATAL, or CRITICAL. The anomaly score
is the corresponding binary value.

Nine-field F1 gives equal weight to seven structural and two anomaly fields.
The three excluded fields are event type, affected entity, and remediation
hint, whose semantic references are not independently validated. Two empty
field values count as a match. Schema validity checks supported field,
type, and allowed-value constraints. Repeated JSON outputs use the first
object. Neither metric proves a correct downstream resolution.

\paragraph{Primary-pair reanalysis.}
The paired comparison uses AR-30B seed-42 predictions and the first 300
DG repeat-1 predictions, aligned to the same prepared test prefix. The
other AR timing seeds reorder this same pool; they do not supply new
quality examples. The frozen references and \texttt{audit\_primary\_pair.py}
reproduce the tables from saved predictions without new inference or
quality-based filtering. The report records source hashes and per-field scores.
A paired percentile bootstrap resamples event-level F1 differences 10,000
times with seed 42. DG minus AR has mean $-0.0181$ and a 95\% interval
$[-0.0289,-0.0086]$. This shows a small DG deficit; no equivalence margin
was selected after observing it. Shared templates and parser references
limit generalization beyond these events.

\paragraph{Timing and serving.}
Both primary models use greedy decoding with repetition penalty 1; DG uses
$K=8$ and a 256-position canvas. Timings pool 900 AR observations across
three runs and 2,000 DG observations on an A100 80\,GB GPU. Median
generation times and their event-bootstrap 95\% intervals are
\arMedian\,ms [\arCiLo, \arCiHi] for AR and
\dgMedian\,ms [\dgCiLo, \dgCiHi] for DG. These intervals are conditional
on the saved runs, not independent hardware replications. Budget counts
include every recorded event, count equality with the budget as on time,
and optionally require schema validity; they exclude queueing, tools,
and semantic outcome checks.

AR metadata records vLLM 0.25.1, torch 2.11.0, and CUDA 13.0. DG's anchor
summary records that stack, but its 2,000-event summary reports the same
timings with vLLM 0.11.0, torch 2.8.0, and CUDA 12.8. The artifacts do
not resolve this conflict, so exact software parity is not established.

\paragraph{Corrections that affect interpretation.}
The current comparison uses raw-log inputs and corrected AR measurements;
earlier gold-record inputs and batched-cost-as-latency claims were invalid.
Earlier smaller-model results also required corrections to decoding settings,
JSON parsing, and field-presence checks. A parser baseline sharing reference
construction was withdrawn. These historical values are not evidence for
the primary pair. The repository retains the detailed audit trail and raw
artifacts; numeric verification does not certify the experimental protocol.

\clearpage
\section{Supporting measurements and ablations}
\label{app:secondary}

\paragraph{Canvas size and refinement effort.}
A joint DG sweep varies canvas sizes $\{64,128,256,512\}$ and refinement
passes $K\in\{2,4,8,16,32\}$, with 200 events per cell. The two
$K=8$ settings below illustrate the tuning opportunity. These are
within-sweep point estimates, separate from the primary-pair reanalysis.

\begin{center}
\begin{tabular}{@{}lrr@{}}
\toprule
DG setting & Median latency (ms) & Nine-field F1 \\
\midrule
Canvas 128, $K=8$ & 368 & 0.864 \\
Canvas 256, $K=8$ & 376 & 0.853 \\
\bottomrule
\end{tabular}
\end{center}

Canvas 128 improves both point estimates in this sample, but does not
establish a global optimum. The paper retains canvas 256 for continuity
with its reference runs. Selecting a replacement needs held-out quality,
validity, and deadline coverage. Reducing $K$ and truncating the output
budget are different interventions: the former reduces refinement effort;
the latter can prevent completion of required fields. Neither establishes
safe interruption of an unfinished canvas.

\paragraph{Output budget versus actual length.}
The AR serial sweep stops near 126 generated tokens once the budget is
sufficient. Budgets 256, 512, and 1024 therefore all yield a mean service
time near 945\,ms. Charging DG for those full canvas budgets instead predicts
379, 498, and 737\,ms. This comparison changes permitted work, not required
answer length. Extrapolating the cost laws to naturally longer responses
requires tasks that actually produce those responses and remain correct.

\paragraph{Serving controls and resource interpretation.}
The AR-30B sweep reaches a \arBatchMult$\times$ throughput multiplier
from batch one to its best tested batch. A separate dense AR-3B control
retains a \cacheOffMult$\times$ multiplier with prefix caching disabled,
compared with \cacheOnMult$\times$ with caching, so shared prefixes alone
do not explain that smaller model's gain. LLaDA-8B, a separate diffusion
baseline, improves throughput by about $1.2\times$ in its tested engine.
These are engine-specific controls, not a universal diffusion batching limit.
Throughput, per-request latency, and average resource time per request must
be kept distinct. Lower serial latency does not by itself reduce the cost
of an always-on GPU; utilization and provisioned capacity also matter.

\paragraph{Scope of the serial-call illustration.}
Figure~\ref{fig:chains} uses DG and dense Qwen2.5-3B, with 100 sequences
at each depth and one call in flight at a time. Prompts are prepared in
advance; answers do not select subsequent prompts. The cumulative median
gap reaches 8.40 seconds at 20 calls while the ratio stays near $2\times$.
This isolates serial generation time, without tool execution, adaptive
decisions, or a measured resolution rate. It cannot establish AR-30B
investigation speedups. With only 100 sequences per depth, the saved p99.9
values are near-maximum observations and are not reliable production targets.

\clearpage
\section{Code-generation limits}
\label{app:code_limits}

\paragraph{Code generation provides a quality counterexample.}
HumanEval~\cite{humaneval} contains 164 Python function-completion problems.
The exploratory single-pass comparison below is serial. Pass@1 is the
fraction of first completions passing the harness tests; syntax validity
is acceptance by Python's parser.

\begin{center}
\begin{tabular}{@{}lrrr@{}}
\toprule
Model & Pass@1 & Syntax valid (\%) & Median latency (ms) \\
\midrule
AR-30B & 0.720 & 93.9 & 800 \\
DG, $K=8$ & 0.287 & 31.1 & 280 \\
\bottomrule
\end{tabular}
\end{center}

DG generates faster here, but its much lower quality prevents a
quality-matched code-generation claim. Finishing a canvas does not
guarantee a syntactically valid or correct program. These results do not
identify whether training, decoding, or another implementation difference
causes the gap.

Other saved code experiments do not close this gap. A nominal deadline
was converted into an AR token budget without enforcing elapsed time;
that is not a measurement of correct completion before a deadline.
A generate--test--repair experiment records DG first-round success of
0.854, conflicting with the single-pass 0.287. Its generated code was
not retained for full audit, and its chain timings exclude test execution.
Those summaries are excluded from the paper's conclusions pending
protocol reconciliation.
